\documentclass[popl.tex]{subfiles}
\begin{document}
\clearpage\section{Examples}

In this section, we will consider three examples of intersections with finite languages. First, parsing can be viewed as a special case of intersection with a singleton language. Second, we will introduce \textit{completion} as an intersection that admits terminal wildcards in fixed locations. Thirdly, we consider syntax repair, where we will intersect a language representing all possible edit paths within a certain distance to determine the location(s) and fill them with the appropriate terminal(s).

\subsection{Recognition as intersection}

In the case of ordinary CFL recognition, the automaton forms a single row and accepts one word:

\begin{figure}[H]
\resizebox{0.5\textwidth}{!}{
  \begin{tikzpicture}[>=stealth', node distance=2.5cm, initial text=$ $]
    \node[state, initial]         (00) {$q_{0,0}$};
    \node[state, right of=00]     (10) {$q_{1,0}$};
    \node[state, right of=10, draw=none]     (20) {$\ldots$};
    \node[state, accepting, right of=20] (30) {$q_{n,0}$};

    \draw [->] (00) edge[below] node{$\sigma_1$} (10);
    \draw [->] (10) edge[below] node{$\sigma_2$} (20);
    \draw [->] (20) edge[below] node{$\sigma_n$} (30);
  \end{tikzpicture}
}
\end{figure}

Since the word is predetermined, we just need to keep track of nonterminal subsets for each substring. So, given a CFG, $G' : \mathcal{G}$ in Chomsky Normal Form (CNF), we can construct a recognizer for strings $\sigma: \Sigma^n$ as follows. Let $2^V$ be our domain, $0$ be $\varnothing$, $\oplus$ be $\cup$, and $\otimes$ be defined as:\vspace{-10pt}

\begin{align}
  X \otimes Z = \big\{\;w \mid \langle x, z\rangle \in X \times Z, (w\rightarrow xz) \in P\;\big\}
\end{align}

\noindent If we define $\hat\sigma_r = \{w \mid (w \rightarrow \sigma_r) \in P\}$, then construct a matrix with nonterminals on the superdiagonal representing each token, $M_0[r+1=c](G', \sigma) = \;\hat\sigma_r$, the fixpoint $M_{i+1} = M_i + M_i^2$ is uniquely determined by the superdiagonal entries. Omitting the exponentiation-by-squaring detail, the ordinary fixedpoint iteration simply fills successive diagonals:\vspace{-10pt}

\begin{equation}
\begin{aligned}
  \hspace{-0.2cm}M_0=
  \begin{pNiceMatrix}[nullify-dots,xdots/line-style=loosely dotted]
    \varnothing & \hat\sigma_1 & \varnothing & \Cdots & \varnothing  \\
    \Vdots      & \Ddots       & \Ddots      & \Ddots & \Vdots       \\
                &              &             &        & \varnothing  \\
                &              &             &        & \hat\sigma_n \\
    \varnothing & \Cdots       &             &        & \varnothing
  \end{pNiceMatrix} & \,,\, M_1=
  \begin{pNiceMatrix}[nullify-dots,xdots/line-style=loosely dotted]
    \varnothing & \hat\sigma_1 & \Lambda     & \Cdots & \varnothing  \\
    \Vdots      & \Ddots       & \Ddots      & \Ddots & \Vdots       \\
                &              &             &        & \Lambda      \\
                &              &             &        & \hat\sigma_n \\
    \varnothing & \Cdots       &             &        & \varnothing
  \end{pNiceMatrix}\,,\,\ldots & \,,\,M_\infty =
  \begin{pNiceMatrix}[nullify-dots,xdots/line-style=loosely dotted]
    \varnothing & \hat\sigma_1 & \Lambda     & \Cdots & \Lambda^*_\sigma \\
    \Vdots      & \Ddots       & \Ddots      & \Ddots & \Vdots           \\
                &              &             &        & \Lambda          \\
                &              &             &        & \hat\sigma_n     \\
    \varnothing & \Cdots       &             &        & \varnothing
  \end{pNiceMatrix}
\end{aligned}
\end{equation}

Once the fixpoint $M_\infty$ is attained, the proposition $[S \in \Lambda^*_\sigma]$~\footnote{Hereinafter, we use Iverson brackets to denote the indicator function of a predicate with free variables, i.e., $[P] \Leftrightarrow \mathds{1}(P)$.} decides language membership, i.e., $[\sigma \in \mathcal{L}(G)]$. So far, this procedure is essentially the textbook CYK algorithm in a linear algebraic notation~\cite{goodman1999semiring} and a well-established technique in the parsing literature~\cite{grune2008}.

\subsection{Completion as intersection}

Let us now consider a problem of intermediate difficulty, wherein we are given a string template admitting edits at fixed locations, which can be filled by any terminal. When intersected with a CFL, this specifies a finite language whose contents are the set of all words consistent with the template. This problem we call \textit{completion}. Formally,

\begin{definition}[Completion]
  Let $\underline\Sigma = \Sigma \cup \{\_\}$, where $\_$ denotes a hole. We denote $\sqsubseteq: \Sigma^n \times \underline\Sigma^n$ as the relation $\{\langle\sigma', \sigma\rangle \mid \sigma_i \in \Sigma \implies \sigma_i' = \sigma_i\}$ and the set of all members $\{\sigma': \Sigma^+ \mid \sigma' \sqsubseteq \sigma\}$ as $\text{H}(\sigma)$. Given a \textit{porous string}, $\sigma: \underline\Sigma^+$ we seek all syntactically valid members, i.e., $A(\sigma)=\text{H}(\sigma)\cap\ell$.
\end{definition}

Here, the FSA takes a similar shape but can have multiple arcs between adjacent states, e.g.:

\begin{figure}[H]
  \resizebox{0.5\textwidth}{!}{
    \begin{tikzpicture}[>=stealth', node distance=2.5cm, initial text=$ $]
      \node[state, initial]                (00) {$q_{0,0}$};
      \node[state, right of=00]            (10) {$q_{1,0}$};
      \node[state, right of=10]            (20) {$q_{2,0}$};
      \node[state, accepting, right of=20] (30) {$q_{3,0}$};

      \draw [->] (00) edge[below]             node{$\sigma_1$} (10);
      \draw [->] (10) edge[below]             node{$\ldots$}   (20);
      \draw [->] (10) edge[below, bend left]  node{$\Sigma_1$} (20);
      \draw [->] (10) edge[below, bend right] node{$\Sigma_{|\Sigma|}$} (20);
      \draw [->] (20) edge[below]             node{$\ldots$}   (30);
      \draw [->] (20) edge[below, bend left]  node{$\Sigma_1$} (30);
      \draw [->] (20) edge[below, bend right] node{$\Sigma_{|\Sigma|}$} (30);
    \end{tikzpicture}
  }
\end{figure}

\noindent This corresponds to a template with two holes, $\sigma = 1$ \_ \_. Suppose the context-free grammar is $G=\{S\rightarrow N O N, O \rightarrow + \mid \times, N \rightarrow 0 \mid 1\}$. This grammar will first be rewritten into CNF as $G'= \{S \rightarrow N L, N \rightarrow 0 \mid 1, O \rightarrow \times \mid +, L \rightarrow O N\}$. Using the powerset algebra we just defined, the matrix fixpoint $M' = M + M^2$ can be computed as follows, shown in the leftmost column below:\vspace{0.3cm}

\input{figures/domain_fixpoints}

\vspace{8pt}The same procedure can be translated, without loss of generality, into the bit domain ($\mathbb{Z}_2^{|V|}$) using a lexicographic nonterminal ordering, however $M_\infty$ in both $2^V$ and $\mathbb{Z}_2^{|V|}$ represents a decision procedure, i.e., $\big[S\in M_\infty[0, 3]\big]\Leftrightarrow \big[M_\infty[0, 3, 3]=\bs\big] \Leftrightarrow \big[A(\sigma) \neq \varnothing\big]$. Since $M_\infty[0, 3] = \{S\}$, we know there is at least one $\sigma' \in A(\sigma)$, but neither $M_\infty$ in $2^V$ or $\mathbb{Z}_2^V$ lets us recover a witness.

%$\{\text{xor}, \land, \top\}$ is a functionally complete set is equivalent to $\mathbb{Z}_2$ $\top := 1, \land := \times, \text{xor} := +$. We can define $=$ as $(a = b) \Leftrightarrow (a \text{ xor } b) \text{ xor } \top \Leftrightarrow (a + b) + \top$.

To witness $\sigma' \in A(\sigma)$, we can translate the matrix exponential to the GRE domain. We first define $X \otimes Z = [X_2 \cdot Z_1, \varnothing, \varnothing, X_1 \cdot Z_0]$ and $X \oplus Z = [X_i \lor Z_i]_{i \in [0, |V|)}$, mirroring $\oplus, \otimes$ from the powerset domain. Since the unit nonterminals $O, N$ can only occur on the superdiagonal, they may be safely ignored by $\otimes$. To solve for $M_\infty$, we proceed by first computing $E_{0, 2}, E_{1, 3}$:\vspace{-8pt}

\begin{small}
\begin{align*}
  E_{0, 2} &= E_{0, j} \cdot E_{j, 2} = E_{0, 1} \otimes E_{1, 2}                         &  E_{1, 3} &= E_{1, j} \cdot E_{j, 3} = E_{1, 2} \otimes E_{2, 3}\\
  &= [L \in E_{0, 2}, \varnothing, \varnothing, S \in E_{0, 2}]                                           &  &= [L \in E_{1, 3}, \varnothing, \varnothing, S \in E_{1, 3}]\\
  &= [O \in E_{0, 1} \cdot N \in E_{1, 2}, \varnothing, \varnothing, N \in E_{0, 1} \cdot L \in E_{1, 2}] &  &= [O \in E_{1, 2} \cdot N \in E_{2, 3}, \varnothing, \varnothing, N \in E_{1, 2} \cdot L \in E_{2, 3}]\\
  &= [E_{0, 1, 2} \cdot E_{1, 2, 1}, \varnothing, \varnothing, E_{0, 1, 1} \cdot E_{1, 2, 0}]             &  &= [E_{1, 2, 2} \cdot E_{2, 3, 1}, \varnothing, \varnothing, E_{1, 2, 1} \cdot E_{2, 3, 0}]
\end{align*}
\end{small}\vspace{-8pt}

\noindent Now we solve for the corner entry $E_{0, 3}$ by dotting the first row and last column, which yields:\vspace{-8pt}

\begin{align*}
  E_{0, 3} &= E_{0, j} \cdot E_{j, 3} = (E_{0, 1} \otimes E_{1, 3}) \oplus (E_{0, 2} \otimes E_{2, 3})\\
%  &= [E_{0, 1, 2} \cdot E_{1, 3, 1}, \varnothing, \varnothing, E_{0, 1, 1} \cdot E_{1, 3, 0}] + [E_{0, 2, 2} \cdot E_{2, 3, 1}, \varnothing, \varnothing, E_{0, 2, 1} \cdot E_{2, 3, 0}]\\
  &= [E_{0, 1, 2} \cdot E_{1, 3, 1} \lor E_{0, 2, 2} \cdot E_{2, 3, 1}, \varnothing, \varnothing, E_{0, 1, 1} \cdot E_{1, 3, 0} \lor E_{0, 2, 1} \cdot E_{2, 3, 0}]
\end{align*}

\noindent Since we only care about $E_{0, 3, 3} \Leftrightarrow [S \in E_{0, 3}]$, we can ignore the first three entries and solve for:\vspace{-8pt}

\begin{align*}
  E_{0, 3, 3} &= E_{0, 1, 1} \cdot E_{1, 3, 0} \lor E_{0, 2, 1} \cdot E_{2, 3, 0}\\
  &= E_{0, 1, 1} \cdot (E_{1, 2, 2} \cdot E_{2, 3, 1}) \lor E_{0, 2, 1} \cdot \varnothing\\
  &= E_{0, 1, 1} \cdot E_{1, 2, 2} \cdot E_{2, 3, 1} \big(= [N \in E_{0, 1}] \cdot [O \in E_{1, 2}] \cdot [N \in E_{2, 3}]\big)\\
  &= 1 \cdot \{+, \times\} \cdot \{0, 1\}
\end{align*}

\noindent Finally, to recover a witness, we can simply $\texttt{choose}\big(1 \cdot \{+, \times\} \cdot \{0, 1\}\big)$  (see Eq.~\ref{eq:choose}).

\subsection{Repair as intersection}\label{sec:repair_ex}

Now, we are ready to consider the general case of syntax repair, in which case the edit locations are not localized but can occur anywhere inside the snippet. In this case, we construct a lattice of all possible edit paths up to a fixed distance. This structure is called a Levenshtein automaton.

\begin{wrapfigure}{r}{0.5\textwidth}
  \vspace{-0.3cm}
  \begin{center}
      \resizebox{0.5\textwidth}{!}{
  \begin{tikzpicture}[
%->, % makes the edges directed
  >=stealth',
  node distance=2.5cm, % specifies the minimum distance between two nodes. Change if necessary.
%  every state/.style={thick, fill=gray!10}, % sets the properties for each ’state’ node
  initial text=$ $, % sets the text that appears on the start arrow
  ]
  \node[state, initial]                (00) {$q_{0,0}$};
  \node[state, right of=00]            (10) {$q_{1,0}$};
  \node[accepting, state, right of=10] (20) {$q_{2,0}$};
  \node[accepting, state, right of=20] (30) {$q_{3,0}$};
  \node[accepting, state, right of=30] (40) {$q_{4,0}$};
  \node[accepting, state, right of=40] (50) {$q_{5,0}$};

  \node[state, above of=00, shift={(-2cm,0cm)}] (01) {$q_{0,1}$};
  \node[state, right of=01]                          (11) {$q_{1,1}$};
  \node[state, right of=11]                          (21) {$q_{2,1}$};
  \node[accepting, state, right of=21]               (31) {$q_{3,1}$};
  \node[accepting, state, right of=31]               (41) {$q_{4,1}$};
  \node[accepting, state, right of=41]               (51) {$q_{5,1}$};

\node[state, above of=01, shift={(-2cm,0cm)}] (0j) {$q_{0,2}$};
\node[state, right of=0j]                          (1j) {$q_{1,2}$};
\node[state, right of=1j]                          (2j) {$q_{2,2}$};
\node[state, right of=2j]                          (3j) {$q_{3,2}$};
\node[accepting, state, right of=3j]               (4j) {$q_{4,2}$};
\node[accepting, state, right of=4j]               (5j) {$q_{5,2}$};

\node[state, above of=0j, shift={(-2cm,0cm)}] (0k) {$q_{0,3}$};
\node[state, right of=0k]                         (1k) {$q_{1,3}$};
\node[state, right of=1k]                         (2k) {$q_{2,3}$};
\node[state, right of=2k]                         (3k) {$q_{3,3}$};
\node[state, right of=3k]                         (4k) {$q_{4,3}$};
\node[accepting, state, right of=4k]              (5k) {$q_{5,3}$};

\draw [->] (00) edge[below] node{$\sigma_1$} (10);
\draw [->] (10) edge[below] node{$\sigma_2$} (20);
\draw [->] (20) edge[below] node{$\sigma_3$} (30);
\draw [->] (30) edge[below] node{$\sigma_4$} (40);
\draw [->] (40) edge[below] node{$\sigma_5$} (50);

\draw [->] (01) edge[below] node{$\sigma_1$} (11);
\draw [->] (11) edge[below] node[shift={(-0.2cm,0cm)}]{$\sigma_2$} (21);
\draw [->] (21) edge[below] node[shift={(-0.2cm,0cm)}]{$\sigma_3$} (31);
\draw [->] (31) edge[below] node[shift={(-0.2cm,0cm)}]{$\sigma_4$} (41);
\draw [->] (41) edge[below] node{$\sigma_5$} (51);

\draw [->] (0j) edge[below] node{$\sigma_1$} (1j);
\draw [->] (1j) edge[below] node{$\sigma_2$} (2j);
\draw [->] (2j) edge[below] node{$\sigma_3$} (3j);
\draw [->] (3j) edge[below] node{$\sigma_4$} (4j);
\draw [->] (4j) edge[below] node{$\sigma_5$} (5j);

\draw [->] (0k) edge[below] node{$\sigma_1$} (1k);
\draw [->] (1k) edge[below] node{$\sigma_2$} (2k);
\draw [->] (2k) edge[below] node{$\sigma_3$} (3k);
\draw [->] (3k) edge[below] node{$\sigma_4$} (4k);
\draw [->] (4k) edge[below] node{$\sigma_5$} (5k);

\draw [->, orange] (00) edge[left] node{$\phantom{\cdot}$} (11);
\draw [->, orange] (10) edge[left] node{$\phantom{\cdot}$} (21);
\draw [->, orange] (20) edge[left] node{$\phantom{\cdot}$} (31);
\draw [->, orange] (30) edge[left] node{$\phantom{\cdot}$} (41);
\draw [->, orange] (40) edge[left] node{$\phantom{\cdot}$} (51);

% Super-knight arcs
\draw [->, red!70] (00) edge[bend right=8] node[anchor=east, shift={(0.1cm,-0.7cm)}]{$\color{red!70}\sigma_3$}         (3j);
\draw [->, red!70] (10) edge[bend right=8] node[anchor=east, shift={(0.1cm,-0.7cm)}]{$\color{red!70}\sigma_4$}         (4j);
\draw [->, red!70] (20) edge[bend right=8] node[anchor=east, shift={(0.1cm,-0.7cm)}]{$\color{red!70}\sigma_5$}         (5j);

\draw [->, red!70] (01) edge[bend left=8] node[anchor=east, shift={(0.1cm,-0.7cm)}]{$\color{red!70}\sigma_3$}         (3k);
\draw [->, red!70] (11) edge[bend left=8] node[anchor=east, shift={(0.1cm,-0.7cm)}]{$\color{red!70}\sigma_4$}         (4k);
\draw [->, red!70] (21) edge[bend left=8] node[anchor=east, shift={(0.1cm,-0.7cm)}]{$\color{red!70}\sigma_5$}         (5k);

\draw [->, red!60] (00) edge node[anchor=east, shift={(0.1cm,-0.8cm)}]{$\color{red!60}\sigma_4$}  (4k);
\draw [->, red!60] (10) edge node[anchor=east, shift={(0.1cm,-0.8cm)}]{$\color{red!60}\sigma_5$}  (5k);

\draw [->, orange] (01) edge[left] node{$\phantom{\cdot}$} (1j);
\draw [->, orange] (11) edge[left] node{$\phantom{\cdot}$} (2j);
\draw [->, orange] (21) edge[left] node{$\phantom{\cdot}$} (3j);
\draw [->, orange] (31) edge[left] node{$\phantom{\cdot}$} (4j);
\draw [->, orange] (41) edge[left] node{$\phantom{\cdot}$} (5j);

\draw [->, orange] (0j) edge[left] node{$\phantom{\cdot}$} (1k);
\draw [->, orange] (1j) edge[left] node{$\phantom{\cdot}$} (2k);
\draw [->, orange] (2j) edge[left] node{$\phantom{\cdot}$} (3k);
\draw [->, orange] (3j) edge[left] node{$\phantom{\cdot}$} (4k);
\draw [->, orange] (4j) edge[left] node{$\phantom{\cdot}$} (5k);

\draw [->, green] (00) edge[bend left=10, left] node{$\phantom{\cdot}$} (01);
\draw [->, green] (10) edge[bend left=10, left] node{$\phantom{\cdot}$} (11);
\draw [->, green] (20) edge[bend left=10, left] node{$\phantom{\cdot}$} (21);
\draw [->, green] (30) edge[bend left=10, left] node{$\phantom{\cdot}$} (31);
\draw [->, green] (40) edge[bend left=10, left] node{$\phantom{\cdot}$} (41);
\draw [->, green] (50) edge[bend left=10, left] node{$\phantom{\cdot}$} (51);

\draw [->, green] (01) edge[bend left=10, left] node{$\phantom{\cdot}$} (0j);
\draw [->, green] (11) edge[bend left=10, left] node{$\phantom{\cdot}$} (1j);
\draw [->, green] (21) edge[bend left=10, left] node{$\phantom{\cdot}$} (2j);
\draw [->, green] (31) edge[bend left=10, left] node{$\phantom{\cdot}$} (3j);
\draw [->, green] (41) edge[bend left=10, left] node{$\phantom{\cdot}$} (4j);
\draw [->, green] (51) edge[bend left=10, left] node{$\phantom{\cdot}$} (5j);

\draw [->, green] (0j) edge[bend left=10, left] node{$\phantom{\cdot}$} (0k);
\draw [->, green] (1j) edge[bend left=10, left] node{$\phantom{\cdot}$} (1k);
\draw [->, green] (2j) edge[bend left=10, left] node{$\phantom{\cdot}$} (2k);
\draw [->, green] (3j) edge[bend left=10, left] node{$\phantom{\cdot}$} (3k);
\draw [->, green] (4j) edge[bend left=10, left] node{$\phantom{\cdot}$} (4k);
\draw [->, green] (5j) edge[bend left=10, left] node{$\phantom{\cdot}$} (5k);

\draw [->, red!85] (00) edge[bend right=11,below] node[shift={(0.5cm,0.3cm)}]{$\color{red}\sigma_2$}    (21);
\draw [->, red!85] (10) edge[bend right=11,below] node[shift={(0.5cm,0.3cm)}]{$\color{red}\sigma_3$}    (31);
\draw [->, red!85] (20) edge[bend right=11,below] node[shift={(0.5cm,0.3cm)}]{$\color{red}\sigma_4$}    (41);
\draw [->, red!85] (30) edge[bend right=11,below] node[shift={(0.5cm,0.3cm)}]{$\color{red}\sigma_5$}    (51);

\draw [->, red!85] (01) edge[bend right=3,below] node[shift={(0.3cm,0.2cm)}]{$\color{red}\sigma_2$}    (2j);
\draw [->, red!85] (11) edge[bend right=3,below] node[shift={(0.3cm,0.2cm)}]{$\color{red}\sigma_3$}    (3j);
\draw [->, red!85] (21) edge[bend right=3,below] node[shift={(0.3cm,0.2cm)}]{$\color{red}\sigma_4$}    (4j);
\draw [->, red!85] (31) edge[bend right=3,below] node[shift={(0.3cm,0.2cm)}]{$\color{red}\sigma_4$}    (5j);

\draw [->, red!85] (0j) edge[bend left=8,below] node[shift={(-0.45cm,-0.55cm)}]{$\color{red}\sigma_2$}    (2k);
\draw [->, red!85] (1j) edge[bend left=8,below] node[shift={(-0.45cm,-0.55cm)}]{$\color{red}\sigma_3$}    (3k);
\draw [->, red!85] (2j) edge[bend left=8,below] node[shift={(-0.45cm,-0.55cm)}]{$\color{red}\sigma_4$}    (4k);
\draw [->, red!85] (3j) edge[bend left=8,below] node[shift={(-0.45cm,-0.55cm)}]{$\color{red}\sigma_5$}    (5k);

%https://tex.stackexchange.com/a/20986/139648
\draw [decorate,decoration={brace,amplitude=10pt,raise=10pt,mirror}] (00.south west) -- (50.south east) node[midway,yshift=-3em]{\textbf{Snippet length}};
\draw [decorate,decoration={brace,amplitude=10pt,raise=20pt}] (00.south west) -- (0k.north west) node[midway,xshift=-1cm,yshift=-1cm,rotate=-54]{\textbf{Edit distance}};
\end{tikzpicture}
}

  \end{center}
  \caption{Levenshtein NFA recognizing $\mathcal{L}\big(L(\sigma: \Sigma^5, 3)\big)$.}\label{fig:lev_nfa}
  \vspace{-0.5cm}
\end{wrapfigure}

As the original construction defined by Schultz and Mihov~\cite{schulz2002fast} contains cycles and $\varepsilon$-transitions, we propose a variant which is $\varepsilon$-free and acyclic. Furthermore, we adopt a symbolic form that supports infinite alphabets and simplifies the description to follow. Illustrated in Fig.~\ref{fig:lev_nfa} is an example of a small Levenshtein automaton recognizing $\mathcal{L}\big(L(\sigma: \Sigma^5, 3)\big)$. Unlabeled arcs accept any terminal from the alphabet, $\Sigma$. Equivalently, this transition system can be viewed as a kind of proof system within an unlabeled lattice. The following construction is equivalent to Schultz and Mihov's original Levenshtein automaton, but is more amenable to our purposes as it does not contain any $\varepsilon$-arcs, and instead uses skip connections to recognize consecutive deletions of varying lengths.

\input{figures/arc_rules}

Note that the same patch can have multiple Levenshtein alignments. $\textsc{Done}$ constructs the final states, which are all states accepting strings $\sigma'$ whose Levenshtein distance $\Delta(\sigma, \sigma') \leq d_{\max}$.

To avoid creating a parallel bundle of arcs for each insertion and substitution point, we instead decorate each arc with a symbolic predicate, accepting or rejecting $\sigma_i$. To distinguish this symbolic variant from the original construction, we highlight the modified rules in orange below.

\begin{prooftree}
  \AxiomC{$i \in [0, n] \phantom{\land} j \in [1, d_{\max}]$}
  \RightLabel{$\duparrow$}
  \UnaryInfC{$(q_{i, j-1} \overset{{\color{orange}[\neq \sigma_{i+1}]}}{\rightarrow} q_{i,j}) \in \delta$}
  \DisplayProof
  \hskip 1.5em
  \AxiomC{$i \in [1, n] \phantom{\land} j \in [1, d_{\max}]$}
  \RightLabel{$\ddiagarrow$}
  \UnaryInfC{$(q_{i-1, j-1} \overset{{\color{orange}[\neq \sigma_i]}}{\rightarrow} q_{i,j}) \in \delta$}
\end{prooftree}
\begin{prooftree}
  \AxiomC{$i \in [1, n] \phantom{\land} j \in [0, d_{\max}]$}
  \RightLabel{$\drightarrow$}
  \UnaryInfC{$(q_{i-1, j} \overset{{\color{orange}[=\sigma_i]}}{\rightarrow} q_{i,j}) \in \delta$}
  \DisplayProof
  \hskip 1.5em
  \AxiomC{$d \in [1, d_{\max}] \phantom{\land} i \in [d + 1, n] \phantom{\land} j \in [d, d_{\max}]$}
  \RightLabel{$\knightarrow$}
  \UnaryInfC{$(q_{i-d-1, j-d} \overset{{\color{orange}[=\sigma_i]}}{\rightarrow} q_{i,j}) \in \delta$}
\end{prooftree}

Using a symbolic predicate in the construction of the NFA avoids the unnecessary creation of $2(|\Sigma| - 1)\cdot|\sigma|\cdot d_{\max}$ arcs and reduces the representation size of the resulting automaton, but does not affect the underlying semantics. Since negation is only permitted over terminals and never propagates further, complementation can be stricken from the usual GRE axioms (Theorem~\ref{thm:witness_construction}).

\begin{wrapfigure}{r}{0.40\textwidth}
\resizebox{0.4\textwidth}{!}{%
\begin{tikzpicture}[
%->, % makes the edges directed
  >=stealth',
  node distance=2.5cm, % specifies the minimum distance between two nodes. Change if necessary.
%  every state/.style={thick, fill=gray!10}, % sets the properties for each ’state’ node
  initial text=$ $, % sets the text that appears on the start arrow
]
  \node[state, initial]                (00) {$q_{0,0}$};
  \node[state, right of=00]            (10) {$q_{1,0}$};
  \node[accepting, state, right of=10] (20) {$q_{2,0}$};
  \node[accepting, state, right of=20] (30) {$q_{3,0}$};

  \node[state, above of=00, shift={(-2cm,0cm)}] (01) {$q_{0,1}$};
  \node[state, right of=01]                     (11) {$q_{1,1}$};
  \node[state, right of=11]                     (21) {$q_{2,1}$};
  \node[accepting, state, right of=21]          (31) {$q_{3,1}$};

  \draw [->] (00) edge[below] node{\tiny{$[= \texttt{(}]$}} (10);
  \draw [->] (10) edge[below] node{\tiny{$[= \texttt{)}]$}} (20);
  \draw [->] (20) edge[below] node{\tiny{$[= \texttt{)}]$}} (30);

  \draw [->] (01) edge[below] node{\tiny{$[= \texttt{(}]$}}                       (11);
  \draw [->] (11) edge[below] node[shift={(-0.2cm,0cm)}]{\tiny{$[= \texttt{)}]$}} (21);
  \draw [->] (21) edge[below] node[shift={(-0.2cm,0cm)}]{\tiny{$[= \texttt{)}]$}} (31);

  \draw [->,orange] (00) edge[left] node{\tiny{$[\neq \texttt{(}]$}} (11);
  \draw [->,orange] (10) edge[left] node{\tiny{$[\neq \texttt{)}]$}} (21);
  \draw [->,orange] (20) edge[left] node{\tiny{$[\neq \texttt{)}]$}} (31);

  \draw [->,green] (00) edge[bend left=10, left] node{\tiny{$[\neq \texttt{(}]$}} (01);
  \draw [->,green] (10) edge[bend left=10, left] node{\tiny{$[\neq \texttt{)}]$}} (11);
  \draw [->,green] (20) edge[bend left=10, left] node{\tiny{$[\neq \texttt{)}]$}} (21);
  \draw [->,green] (30) edge[bend left=10, left] node{\tiny{$[=.]$}} (31);


  \draw [->, red!85] (00) edge[bend right=11,below] node[shift={(0.2cm,0.8cm)}]{\tiny{$[= \texttt{)}]$}}    (21);
  \draw [->, red!85] (10) edge[bend right=11,below] node[shift={(0.2cm,0.8cm)}]{\tiny{$[= \texttt{)}]$}}    (31);
\end{tikzpicture}

}
\caption{Simple Levenshtein automaton.}\label{fig:ex_atm}

\vspace{0.3cm}
\resizebox{0.4\textwidth}{!}{%
\input{figures/partial_order}
}
\caption{Pairing function over $L(\sigma: \Sigma^3, 1)$.}\label{fig:pairing_fun}

\vspace{0.3cm}
\begin{center}
\resizebox{0.35\textwidth}{!}{%
\input{figures/adj_mat}
}
\end{center}
\vspace{-0.3cm}
\caption{Adjacency and reachability matrix.}\label{fig:reach_matr}

\vspace{0.3cm}
\resizebox{0.4\textwidth}{!}{
\input{figures/pc_init}
}
\caption{Initial parse chart configuration.}\label{fig:initial_pc}

\vspace{0.3cm}
\resizebox{0.4\textwidth}{!}{
\input{figures/pc_final}
}

\caption{Final parse chart configuration.}\label{fig:final_pc}

\begin{center}
\resizebox{0.4\textwidth}{!}{
  \includegraphics{figures/gre}
}
\end{center}
\vspace{-0.3cm}
\caption{Regular expression denoting $\mathcal{L}(G_\cap)$.}\label{fig:re_tree}
\end{wrapfigure}

As a concrete example, suppose we have the string, $\err\sigma=\texttt{())}$ and wish to balance the parentheses. We will initially have the Levenshtein automaton, $A$, depicted in Fig.~\ref{fig:ex_atm}. To check for non-emptiness, we will perform the following procedure. Suppose we have a CNF CFG, $G'= \big\{S \rightarrow L R, S \rightarrow L F, S \rightarrow S S, F \rightarrow S R, L \rightarrow \hspace{-0.05cm}\texttt{(}, R \rightarrow\hspace{-0.05cm}\texttt{)}\big\}$ and let us assume an ordering of $S, F, L, R$ on $V$.

First, we need to order the automata states by increasing longest-path distance from $q_0$. One approach would be to topologically sort the adjacency matrix. While some form of sorting is unavoidable for arbitrary ANFAs, if we know ahead of time that our structure is a Levenshtein automaton, we can simply enumerate its state space by increasing Manhattan distance from the origin. % using, e.g., the Cantor pairing function to construct a valid ordering. This ordering will form the row and column indices of our intersection matrix, and each entry will represent the existence of some path between a two states yielding a given nonterminal.
So, a valid ordering on $Q$ would be $q_{00}, q_{01}, q_{10}, q_{11}, q_{20}, q_{21}, q_{30}, q_{31}$. Now, we want to compute whether $[\mathcal{L}(G')\cap \mathcal{L}(A) \neq \varnothing]$.

Under such an ordering, the adjacency matrix takes an upper triangular form and becomes the template for the initial parse chart, $M_0$ (Fig.~\ref{fig:initial_pc}). Each entry of this chart corresponds to a vector of expressions $E^{|V|}$ with at least one expression denoting a nonempty language. Likewise, the reachability matrix signifies a subset of state pairs which can participate in the language intersection. The adjacency and reachability matrices will always cover the expression vectors of the initial and final parse charts, respectively. In other words, we may safely ignore absent $\langle q, q'\rangle$ pairs in the reachability matrix, as these state pairs definitely cannot participate in the intersection.

From the reachability matrix we can construct the parse chart via matrix exponentiation. We note that n-step reachability constrains n-step parseability, i.e., $\sum_{i=0}^n A^i[q, q'] = \ws \vdash M_n[q, q', v] = \ws$, thus we can avoid substantial work via memoization. In this example, since $M_\infty[q_{00}, q_{31}, S] = \bs$, this implies that $\mathcal{L}(A)\cap \mathcal{L}(G') \neq \varnothing$, hence $\text{LED}(\sigma, G) = 1$. Using the same matrix, we will then perform a second pass to construct regular expressions representing finite languages for each nonempty constituent. Once again, we can skip $\langle q, q', v\rangle$ entries when $M_\infty[q, q', v] = \ws$ to hasten convergence.

\enlargethispage{4\baselineskip}
Just as before, we will define $\oplus, \otimes$ using the provenance semiring defined in Thm.~\ref{thm:witness_construction}. Finally, we will repeat the matrix exponential, using $M_\infty$ in the binary domain as a guide. This allows us to construct the regular expression tree for $S_\cap = q_{00}Sq_{20}\vee q_{00}Sq_{31}$ shown in Fig.~\ref{fig:re_tree}. Once constructed, decoding becomes simply a matter of invoking \texttt{choose}$(S_\cap)$. In this case, there are only a few choices, but in general, there can be a vast multitude.

\clearpage
\ifSubfilesClassLoaded{\par\clearpage\bibliography{../bib/acmart}}{}
\end{document}
